\documentclass[11pt,a4paper]{article}
\usepackage{gtu-el-yazisi}
\begin{document}
\notbilgi{İşaretler ve Sistemler}{Uygulama notları -- Fourier serisi ve Fourier dönüşümü (sayfa 22--41)}{isaretler-uygulama.pdf}
% \notbilgi already prints the bilingual "transcribed with Claude, may contain errors" box.

\sayfa{22}
\[
x_k = 1\int_0^1 (1-t)\cdot e^{-jk2\pi t}\,dt
 = \underbrace{\int_0^1 e^{-jk2\pi t}\,dt}_{I_1} - \underbrace{\int_0^1 t\cdot e^{-jk2\pi t}\,dt}_{I_2}
 = I_1 - I_2 \ \checkmark
\]
\textcolor{morkalem}{$\displaystyle\int e^{-at}\,dt = -\frac{1}{a}\,e^{-at}$}
\[
I_1 = \int_0^1 e^{-jk2\pi t}\,dt = \frac{-1}{jk2\pi}\,e^{-jk2\pi t}\Big|_0^1
 = \frac{-1}{jk2\pi}\left[e^{-j2\pi k}-1\right] = \frac{1}{j2\pi k}\left(1-e^{-j2\pi k}\right)\ \checkmark
\]
\[
I_2 = \int_0^1 t\cdot e^{-jk2\pi t}\,dt \qquad\qquad
\textcolor{kirmizikalem}{\text{Kısmi integrasyon: }\int u\,dv = u\cdot v - \int v\,du}
\]
\[
\begin{aligned}
u &= t & e^{-jk2\pi t}\,dt &= dv\\
du &= du\text{\footnotemark} & \frac{-1}{jk2\pi}\,e^{-jk2\pi t} &= v
\end{aligned}
\]
\footnotetext{Claude notu: $du = dt$ olmalı; s.~23'te öğrenci tarafından $du = dt$ olarak düzeltilmiş.}
\[
I_2 = t\cdot\frac{-1}{jk2\pi}\,e^{-jk2\pi t}\Big|_0^1 + \int_0^1 \frac{1}{jk2\pi}\,e^{-jk2\pi t}\,dt
\]
\[
= \frac{-1}{jk2\pi}\,e^{-jk2\pi} - 0 + \frac{1}{jk2\pi}\cdot
 \underbrace{\frac{-1}{jk2\pi}\cdot e^{-jk2\pi t}\Big|_0^1}_{\frac{1}{j2\pi k}\left(1-e^{-j2\pi k}\right)}
\]
\[
\kutula{x(t) = \sum_{k=-\infty}^{\infty} \frac{1}{j2\pi k}\cdot e^{-jk2\pi t}}\text{\footnotemark}
\]
\footnotetext{Claude notu: sentez denkleminde muhtemelen $e^{+jk2\pi t}$ olmalı.}

\sayfa{23}
{\small\itshape (Claude notu: bu sayfanın üst kısmı s.~22 ile aynıdır; yalnızca yeni eklenen satırlar aktarıldı. $I_2$ içinde $u$, $v$, $du$, $dv$ kısımları pembe kalemle işaretlenmiş.)}
\newcommand{\iptal}[1]{\tikz[baseline=(gtuIptal.base)]{\node[inner sep=1pt](gtuIptal){$\displaystyle #1$};
  \draw[kursunkalem,line width=0.6pt] (gtuIptal.south west) -- (gtuIptal.north east);}}
\[
= \frac{-1}{j2\pi k}\cdot e^{-j2\pi k} + \frac{1}{j^2(2\pi k)^2}\cdot\left(1-e^{-j2\pi k}\right)
\]
\[
x_k = I_1 - I_2 = \frac{1}{j2\pi k} - \iptal{\frac{1}{j2\pi k}\,e^{-j2\pi k}}
 + \iptal{\frac{1}{j2\pi k}\,e^{-j2\pi k}}
 + \frac{1}{(2\pi k)^2}\Bigl(1-\underbrace{e^{-j2\pi k}}_{=1}\Bigr) = \frac{1}{j2\pi k}
 \ \nearrow\ \text{(kutudaki } x(t)\text{)}
\]

\sayfa{24}
{\small\itshape (Claude notu: s.~23 ile aynı sayfa. Pembe kalemle eklenenler: üstteki ``$= \frac{-1}{j2\pi k}e^{-j2\pi k} + \dots$'' satırı $I_2$ olarak etiketlenip $\checkmark$ ile onaylanmış, sadeleşen iki terimin üstü çizilmiş ve $e^{-j2\pi k}$ daire içine alınmış.)}
\[
\underbrace{e^{-j2\pi k}}_{=1}\ \longrightarrow\ \textcolor{morkalem}{e^{j\theta} = \cos\theta + j\sin\theta}
\]

\sayfa{25}
\begin{ornek}[(Soru)]
Aşağıda verilen periyodik işaretin Fourier serisi açılımını bulunuz.

\begin{center}
\begin{tikzpicture}[x=0.8cm,y=1.05cm]
  \draw[dotted,kursunkalem,line width=0.8pt] (-5.3,1) -- (5.3,1) (-5.3,-1) -- (5.3,-1);
  \foreach \a in {-5,-3,-1,1,3,5} \draw[dashed,kursunkalem] (\a,-1) -- (\a,1);
  \draw[-{Stealth},murekkep] (-5.9,0) -- (6.4,0) node[right] {$t$};
  \draw[-{Stealth},murekkep] (0,-1.5) -- (0,1.6) node[right] {$x(t)$};
  \foreach \n in {-2,-1,1,2} \draw[murekkep,thick] ({2*\n-1},-1) -- ({2*\n+1},1);
  \draw[kirmizikalem,very thick] (-1,-1) -- (1,1);
  \fill[kirmizikalem] (0,0) circle (2.2pt) (2,0) circle (2.2pt);
  \node[below right,kirmizikalem,font=\small,inner sep=1pt] at (0.05,-0.05) {$(0,0)$};
  \node[above left,font=\small,inner sep=1pt] at (-1,0) {$-T$};
  \node[above right,font=\small,inner sep=1pt] at (1,0) {$T$};
  \node[below right,font=\small,inner sep=1pt] at (2,0) {$2T$};
  \node[above left,font=\small,inner sep=1pt] at (3,0) {$3T$};
  \node[below right,font=\small,inner sep=1pt] at (4,0) {$4T$};
  \node[above left,font=\small,inner sep=1pt] at (0,1) {$1$};
  \node[below right,font=\small,inner sep=1pt] at (0,-1) {$-1$};
  \node at (-5.7,0.45) {$\cdots$};
  \node at (5.9,0.45) {$\cdots$};
\end{tikzpicture}
\end{center}
\end{ornek}

\textcolor{yesilkalem}{slope: $\dfrac{2}{2T} = \dfrac{1}{T}$} $\checkmark$ \quad
\textcolor{yesilkalem}{D.D!\quad $\dfrac{1}{T}\cdot x = y$} $\checkmark$ \qquad
\textcolor{kirmizikalem}{$x(t) = \dfrac{t}{T}$}
\[
\underline{T_0 = 2T}\,, \qquad \omega_0 = \frac{2\pi}{T_0} = \frac{2\pi}{2T} = \frac{\pi}{T}\ \checkmark
\]
\[
x_k = \frac{1}{T_0}\int_{T_0} x(t)\cdot e^{-jk\omega_0 t}\,dt \ \checkmark
\]
\[
= \frac{1}{2T}\int_{-T}^{T} \frac{t}{T}\cdot e^{-jk\frac{\pi}{T}\cdot t}\,dt
 = \frac{1}{2T^2}\int_{-T}^{T} t\cdot e^{-jk\frac{\pi}{T}\cdot t}\,dt
 \qquad \textcolor{kirmizikalem}{\text{* Kısmi integrasyon!}}
\]
\[
\begin{aligned}
u &= t & e^{-jk\frac{\pi}{T}\cdot t}\,dt &= dv\\
du &= dt & \frac{-1}{jk\frac{\pi}{T}}\,e^{-jk\frac{\pi}{T}t} &= v
\end{aligned}
\]

\sayfa{26}
\[
= \frac{1}{2T^2}\left\{ t\cdot\frac{-1}{jk\frac{\pi}{T}}\cdot e^{-jk\frac{\pi}{T}t}\Big|_{-T}^{T}
 + \int_{-T}^{T}\frac{1}{jk\frac{\pi}{T}}\,e^{-jk\frac{\pi}{T}t}\,dt \right\}
\]
\[
= \frac{1}{2T^2}\left\{ \frac{-1}{jk\frac{\pi}{T}}\cdot\left[T\cdot e^{-jk\pi} + T\cdot e^{jk\pi}\right]
 + \frac{1}{jk\frac{\pi}{T}}\cdot\frac{-1}{jk\frac{\pi}{T}t}\cdot e^{-jk\frac{\pi}{T}t}\Big|_{-T}^{T} \right\}\qquad\text{\footnotemark}
\]
\footnotetext{Claude notu: paydadaki $t$ muhtemelen fazladan yazılmış; $\frac{-1}{jk\frac{\pi}{T}}$ olmalı (bir sonraki satırda da öyle kullanılmış).}
\[
= \frac{1}{2T^2}\left\{ \frac{jT^2}{\pi k}\left[e^{-j\pi k} + e^{j\pi k}\right]
 + \frac{1}{\left(k\cdot\frac{\pi}{T}\right)^2}\left[e^{-j\pi k} - e^{j\pi k}\right] \right\}
 \qquad
 \textcolor{kirmizikalem}{\begin{aligned} e^{j\pi n} &= (-1)^n\\ e^{-j\pi n} &= (-1)^n \end{aligned}}
\]
\[
= \frac{j}{2\pi k}\cdot 2\cdot(-1)^k + \frac{1}{2(\pi k)^2}\cdot 0 = \frac{j}{\pi k}\cdot(-1)^k = \frac{(-1)^{k+1/2}}{\pi k}
\]
\[
x(t) = \sum_{k=-\infty}^{\infty} \frac{(-1)^{k+1/2}}{\pi k}\,e^{jk\frac{\pi}{T}\cdot t}
\]

\sayfa{27}
{\small\itshape (Claude notu: s.~26 ile aynı sayfa; kırmızı $e^{\pm j\pi n} = (-1)^n$ ifadeleri kutu içine alınmış, $e^{\mp j\pi k}$ ve $(-1)^k$ terimlerinin altı çizilmiş. Pembe kalemle eklenenler:)}
\[
\textcolor{morkalem}{\underline{\underline{e^{-j2\pi k} = 1}}}
\]
\[
\textcolor{morkalem}{\begin{aligned}
&\left(e^{j\pi}\right)^n\\
&e^{j\pi} = \cos\pi + \iptal{j\sin\pi} = -1 \ \longrightarrow\ (-1)^n\\
&e^{-j\pi} = -1 \ \longrightarrow\ \left(e^{-j\pi}\right)^n = (-1)^n
\end{aligned}}
\]

\sayfa{28}
\begin{ornek}[(Soru)]
\begin{minipage}[c]{0.42\linewidth}
\centering
\begin{tikzpicture}[scale=1.25]
  \draw[-{Stealth},murekkep] (-1.6,0) -- (1.7,0) node[right] {$t$};
  \draw[-{Stealth},murekkep] (0,-1.35) -- (0,1.45) node[right] {$x(t)$};
  \draw[murekkep,thick] (-1,0) -- (0,1) -- (0,-1) -- (1,0);
  \fill[yesilkalem] (-1,0) circle (1.3pt) (1,0) circle (1.3pt);
  \node[below,font=\small] at (-1,0) {$-1$};
  \node[above,font=\small] at (1,0) {$1$};
  \node[right,font=\small] at (0,1) {$1$};
  \node[left,font=\small] at (0,-1) {$-1$};
\end{tikzpicture}
\end{minipage}\hfill
\begin{minipage}[c]{0.55\linewidth}
$x(t)$ işaretinin FT'nu bulunuz.
\[
\textcolor{yesilkalem}{X(f) = \int_{-\infty}^{\infty} x(t)\,e^{-j2\pi ft}\,dt}
\]
\end{minipage}

\medskip
\textcolor{yesilkalem}{slope$_1 = +1$ \quad D.D: $1\cdot(x+1) = y$} \qquad\qquad
\textcolor{yesilkalem}{slope$_2 = +1$ \quad D.D: $1(x-1) = y$}
\end{ornek}

\[
X(f) = \int_{-1}^{0} (t+1)\cdot e^{-j2\pi ft}\,dt + \int_{0}^{1} (t-1)\cdot e^{-j2\pi ft}\,dt
\]
\[
= \int_{-1}^{0} t\cdot e^{-j2\pi ft}\,dt + \int_{-1}^{0} e^{-j2\pi ft}\,dt
 + \int_{0}^{1} t\cdot e^{-j2\pi ft}\,dt - \int_{0}^{1} e^{-j2\pi ft}\,dt
\]
\[
= \frac{j}{\pi f}\left[1 - \operatorname{sinc}(2f)\right]
\]

\sayfa{29}
{\small\itshape (Claude notu: s.~28 ile aynı sayfa; FT tanımı kutu içine alınıp $\checkmark$ ile, sonuç da $\checkmark$ ile işaretlenmiş. Pembe kalemle eklenenler:)}

\textcolor{yesilkalem}{D.D: $1\cdot(x+1) = y$} $\textcolor{morkalem}{\ \longrightarrow\ x(t) = t+1}$ \qquad\qquad
\textcolor{yesilkalem}{D.D: $1(x-1) = y$} \ $\textcolor{morkalem}{x(t) = t-1}$
\[
\textcolor{morkalem}{\kutula{\int t\cdot e^{-at}\,dt}}
\qquad\qquad
\textcolor{morkalem}{\kutula{\frac{\sin(\pi x)}{\pi x} = \operatorname{sinc}(x)}}
\]

\sayfa{30}
\begin{ornek}[(Soru)]
$x(t) = \Pi(2t+3) + \Pi(3t-2)$ \ işaretinin Fourier transformunu bulunuz.
\[
\textcolor{morkalem}{x(t) \xrightarrow{\ \mathcal{F}\ } X(f)}
\]
\end{ornek}

\noindent
\begin{minipage}[c]{0.66\linewidth}
\textcolor{yesilkalem}{Hatırlatma:}
\[
\textcolor{yesilkalem}{\begin{aligned}
A\cdot\Pi(t/\tau) &\xrightarrow{\ \textcolor{morkalem}{\mathcal{F}}\ } A\cdot\tau\operatorname{sinc}(f\tau) \ \textcolor{morkalem}{\longleftarrow}\\
x(t-t_0) &\xrightarrow{\ \textcolor{morkalem}{\mathcal{F}}\ } X(f)\,e^{-j2\pi ft_0} \ \textcolor{morkalem}{\longleftarrow}
\end{aligned}}
\]
\end{minipage}\hfill
\begin{minipage}[c]{0.32\linewidth}
\centering
\begin{tikzpicture}[x=1.1cm,y=1cm]
  \draw[-{Stealth},morkalem] (-1.5,0) -- (1.6,0) node[right] {$t$};
  \draw[-{Stealth},morkalem] (0,0) -- (0,1.2);
  \draw[morkalem,thick] (-0.8,0) -- (-0.8,0.85) -- (0.8,0.85) -- (0.8,0);
  \node[morkalem,above right,font=\small,inner sep=1pt] at (0,0.85) {$A$};
  \node[morkalem,below,font=\small] at (-0.8,0) {$-\tau/2$};
  \node[morkalem,below,font=\small] at (0.8,0) {$\tau/2$};
  \node[morkalem,font=\small] at (0.6,1.35) {$A\cdot\Pi(t/\tau)$};
\end{tikzpicture}
\end{minipage}

\[
\Pi\left(\frac{t+3/2}{1/2}\right) \longrightarrow \frac{1}{2}\cdot\operatorname{sinc}(f/2)\cdot e^{-j2\pi f(-3/2)}
\]
\[
\Pi\left(\frac{t-2/3}{1/3}\right) \longrightarrow \frac{1}{3}\operatorname{sinc}\left(\frac{f}{3}\right)e^{-j2\pi f\cdot 2/3}
\]
\[
X(f) = \frac{1}{2}\operatorname{sinc}\left(\frac{f}{2}\right)\cdot e^{+j2\pi f\cdot 3/2}
 + \frac{1}{3}\operatorname{sinc}\left(\frac{f}{3}\right)e^{-j2\pi f\cdot 2/3}
\]

\sayfa{31}
{\small\itshape (Claude notu: s.~30 ile aynı sayfa; $\Pi(2t+3)$ daire içine, ilk dönüşüm sonucu kutu içine alınmış. Pembe kalemle eklenenler:)}
\[
\textcolor{morkalem}{\Pi(2t+3) = }\ \Pi\left(\frac{t+3/2}{1/2}\right)
 \longrightarrow
 \kutula{\frac{1}{2}\cdot\operatorname{sinc}(f/2)\cdot e^{-j2\pi f(-3/2)}}
\]
\[
\textcolor{morkalem}{2(t+3/2) \qquad\quad t_0 = -3/2 \qquad\quad \tau = \frac{1}{2} \qquad\quad f\cdot\tau = \frac{f}{2}}
\]

\sayfa{32}
\begin{ornek}[(Soru)]
$x(t) = \overbrace{4\cdot\Pi(t/4)}^{\textcolor{morkalem}{\hat{x}(t)}}\cdot\cos(2\pi f_0 t)$ \ işaretinin fourier transformunu bulunuz.
\end{ornek}

\textcolor{morkalem}{* Modülasyon}\\
\textcolor{yesilkalem}{Hatırlatma:} \
$\textcolor{yesilkalem}{\underline{x(t)}\cdot\underline{\cos(2\pi f_0 t)} \longrightarrow \frac{1}{2}\left[X(f-f_0) + X(f+f_0)\right]}$ \ \textcolor{morkalem}{$\checkmark$}
\[
\hat{x}(t) = 4\cdot\Pi(t/4) \longrightarrow \hat{X}(f) = 16\cdot\operatorname{sinc}(f\cdot 4)
\]
\[
X(f) = \frac{1}{2}\cdot 16\cdot\operatorname{sinc}\bigl(4(f-f_0)\bigr) + \frac{1}{2}\cdot 16\cdot\operatorname{sinc}\bigl(4(f+f_0)\bigr)
\]
\[
= 8\operatorname{sinc}\bigl(4(f-f_0)\bigr) + 8\operatorname{sinc}\bigl(4(f+f_0)\bigr)
\]

\sayfa{33}
{\small\itshape (Claude notu: s.~32 ile aynı sayfa; sorudaki $\hat{x}(t)$ etiketi bu sayfada tamamlanmış, sonuç $\checkmark$ ile onaylanmış. Pembe kalemle eklenenler:)}
\[
\textcolor{morkalem}{\hookrightarrow\ X(f) = \frac{1}{2}\left[\hat{X}(f-f_0) + \hat{X}(f+f)\right]}\qquad\text{\footnotemark}
\]
\footnotetext{Claude notu: muhtemelen $\hat{X}(f+f_0)$ olmalı.}
\[
\underset{\textcolor{morkalem}{A=4\quad \tau=4}}{\hat{x}(t) = 4\cdot\Pi(t/4)} \longrightarrow
\underline{\underline{\hat{X}(f)}} = \underset{\textcolor{morkalem}{A\tau}}{\underline{\underline{16}}}\cdot\operatorname{sinc}\bigl(f\cdot\underset{\textcolor{morkalem}{\tau}}{4}\bigr)
\]
\[
X(f) = \underbrace{\frac{1}{2}\cdot 16\cdot\operatorname{sinc}\bigl(4(f-f_0)\bigr)}_{\textcolor{morkalem}{X(f-f_0)}}
 + \underbrace{\frac{1}{2}\cdot 16\cdot\operatorname{sinc}\bigl(4(f+f_0)\bigr)}_{\textcolor{morkalem}{X(f+f_0)}}
\]
\[
\textcolor{morkalem}{A\cdot\Pi(t/\tau) \xrightarrow{\qquad} A\tau\operatorname{sinc}(f\tau)}
\]

\sayfa{34}
\begin{ornek}[14)]
$x(t) = \dfrac{1}{1+t^2}$ \qquad Find the Fourier Transform.
\end{ornek}
\[
e^{-\alpha|t|} \longrightarrow \frac{2\alpha}{\alpha^2 + (2\pi f)^2}
 = \frac{2\alpha}{4\pi^2}\cdot\frac{1}{\frac{\alpha^2}{4\pi^2} + f^2}
\]
\[
\left.\begin{aligned}
x(t) &\longrightarrow X(f)\\
X(t) &\longrightarrow x(-f)
\end{aligned}\ \right\}\ \underline{\underline{\text{Duality}}}
\]
\[
X(t) = \frac{1}{1+t^2} \xrightarrow{\ \mathcal{F}\ } x(-f)
\qquad\qquad
\frac{2\alpha}{4\pi^2}\,\frac{1}{\frac{\alpha^2}{4\pi^2} + t^2} \xrightarrow{\ \mathcal{F}\ } e^{-\alpha|-f|} = e^{-\alpha|f|}
\]
\[
\frac{1}{\frac{\alpha^2}{4\pi^2} + t^2} \longrightarrow \frac{4\pi^2}{2\alpha}\cdot e^{-\alpha|f|}
\]
\[
\frac{\alpha^2}{4\pi^2} = 1 \qquad \alpha^2 = 4\pi^2 \ \to\ \underline{\underline{\alpha = 2\pi}}
\]
\[
\frac{1}{1+t^2} \xrightarrow{\ \mathcal{F}\ } \pi\cdot e^{-2\pi|f|}
\]

\sayfa{35}
{\small\itshape (Claude notu: s.~34 ile aynı sayfa; duality çifti yeşil kutu içine alınmış, ``Fourier Transform'' altı çizilmiş. Yeşil/pembe kalemle eklenenler:)}
\[
x(t) = \underline{\underline{\frac{1}{1+t^2}}} \overset{\textcolor{morkalem}{?}}{\textcolor{morkalem}{=}} \textcolor{morkalem}{Y(t)}
\qquad\qquad
\textcolor{yesilkalem}{X(f) = \int_{T_0} x(t)\,e^{\cdots}\,dt}
\]
\[
\underset{\textcolor{morkalem}{y(t)}}{e^{-\alpha|t|}} \xrightarrow{\ \mathcal{F}\ }
\underset{\textcolor{morkalem}{Y(f)}}{\frac{2\alpha}{\alpha^2 + (2\pi f)^2}}
\qquad
\textcolor{yesilkalem}{\hookrightarrow\ \begin{aligned}[t]
Y(t) &= \frac{2\alpha}{4\pi^2}\cdot\frac{1}{\frac{\alpha^2}{4\pi^2} + t^2}\\
&= \frac{1}{\pi}\,\kutula{\frac{1}{1+t^2}}
\end{aligned}
\qquad
\begin{aligned}[t]
\frac{\alpha^2}{4\pi^2} &= 1\\
\alpha^2 &= 4\pi^2\\
\underline{\alpha} &\underline{{}= 2\pi}
\end{aligned}}
\]
\[
\textcolor{yesilkalem}{\kutula{Y(t) = \frac{1}{\pi}\cdot x(t)}\ \downarrow}
\]

\sayfa{36}
{\small\itshape (Claude notu: s.~34'teki sayfanın başka bir açıklamalı hâli; s.~35'teki $Y(t)=\frac{1}{\pi}x(t)$ eklemeleri bu sayfada yok. $e^{-\alpha|t|} \to \frac{2\alpha}{\alpha^2+(2\pi f)^2}$ çifti yeşil kutuya alınmış, iki tarafı pembe daire içinde; $x(-f)$ altı çizilmiş. Pembe kalemle eklenenler:)}
\[
\textcolor{morkalem}{\overset{y(t)}{e^{-\alpha|t|}} \xrightarrow{\ \mathcal{F}\ } \overset{Y(t)}{\frac{2\alpha}{\alpha^2 + (2\pi f)^2}}}
\qquad\qquad
\textcolor{morkalem}{\begin{aligned}[t]
Y(t) &= \frac{2\alpha}{\alpha^2 + (2\pi f)^2}\ \nearrow\ y(-f)\\[1mm]
Y(f) &= \frac{2\alpha}{\alpha^2 + (2\pi f)^2}
\end{aligned}}
\]
\[
\textcolor{morkalem}{\frac{2\alpha}{\alpha^2 + (2\pi t)^2} \xrightarrow{\ \mathcal{F}\ } e^{-\alpha|-f|} = e^{-\alpha|f|}}
\]

\sayfa{37}
{\small\itshape (Claude notu: s.~32--33'teki sayfanın alt kısmı; yatay pembe çizginin altı. Bu bölüm çoğunlukla pembe kalemle yazılmış.)}
\[
\kutula{\textcolor{morkalem}{\overset{\textcolor{yesilkalem}{x(t)}}{A\cdot\Pi(t/\tau)} \xrightarrow{\qquad\quad} \underset{\textcolor{yesilkalem}{X(f)}}{A\tau\operatorname{sinc}(f\tau)}}}
\qquad\qquad
\textcolor{morkalem}{\begin{aligned}
x(t) &\longrightarrow X(f)\\
X(t) &\longrightarrow x(-f)
\end{aligned}}
\]
\begin{center}
\begin{tikzpicture}[x=1cm,y=1cm]
  \draw[-{Stealth},morkalem] (-1.6,0) -- (1.7,0) node[right] {$t$};
  \draw[-{Stealth},morkalem] (0,0) -- (0,1.3);
  \draw[morkalem,thick] (-0.9,0) -- (-0.9,0.9) -- (0.9,0.9) -- (0.9,0);
  \node[morkalem,above right,font=\small,inner sep=1pt] at (0,0.9) {$A$};
  \node[morkalem,below,font=\small] at (-0.9,0) {$-\tau/2$};
  \node[morkalem,below,font=\small] at (0.9,0) {$\tau/2$};
  \begin{scope}[xshift=5.2cm]
    \draw[-{Stealth},morkalem] (-2.3,0) -- (2.4,0);
    \draw[-{Stealth},morkalem] (0,-0.4) -- (0,1.4);
    \draw[morkalem,thick,smooth,domain=-12:12,samples=120] plot ({\x*0.18},{1.15*sin(\x r)/\x});
    \node[yesilkalem,font=\small] at (-0.8,1.2) {$X(f)$};
  \end{scope}
\end{tikzpicture}
\end{center}
\[
\textcolor{morkalem}{\text{Duality!}\qquad \underset{\sim\sim\sim}{A\cdot\tau\cdot\operatorname{sinc}(t\tau)} \xrightarrow{\ \ \mathcal{F}\ \ } A\cdot\Pi(f/\tau)}
\]
\[
\textcolor{yesilkalem}{X(t) \xrightarrow{\qquad\quad} x(-f) = A\cdot\Pi(-f/\tau) = A\cdot\Pi(f/\tau)\ \longleftarrow}
\]

\noindent\sayfa{38}
\sekilnotu[38]{Bu sayfa basılı materyal içeriyor (çoktan seçmeli sınav sorusunun ekran görüntüsü, 5. soru, 5 puan); dijitalleştirilmedi. Öğrenci $1+j$ şıkkını daire içine almış.}
El yazısı eklemeler:
\[
\ldots = \sqrt{2}\left[\cos\left(\frac{9\pi}{4}\right) + j\sin\left(\frac{9\pi}{4}\right)\right] = 1+j
\qquad\qquad
e^{j\theta} = \cos\theta + j\sin\theta
\]

\noindent\sayfa{39}
\sekilnotu[39]{Bu sayfa basılı materyal içeriyor (çoktan seçmeli sınav sorusunun ekran görüntüsü, 2. soru, 8 puan: verilen sistemin doğrusallık / zamanla değişmezlik özellikleri); dijitalleştirilmedi.}
El yazısı eklemeler:
\[
a\cdot x(t) \xrightarrow{\ T.\ } a\cdot y(t)
\]
\[
\begin{aligned}
x(t) &\longrightarrow y(t) = t^2\cdot x(t-1)\\
x_1(t) &\longrightarrow y_1(t) = t^2\cdot x_1(t-1)\\
x_1(t) = a\cdot x(t) &\longrightarrow y_1(t) = \underline{t^2\cdot a\cdot x(t-1)}
\end{aligned}
\]
\[
T\left[a x(t)\right] = a\cdot\underset{\sim\sim\sim}{T\left[x(t)\right]}
 = \underline{a\cdot t^2\cdot x(t-1)}
\]

\sayfa{40}
{\small\itshape (Claude notu: s.~39 ile aynı basılı soru. Öğrenci ``Lineer değildir'' içeren iki şıkkı karalamış, son şıkkı (\textit{Lineer, zamanla değişir}) daire içine alıp ``Lineer'' ve ``zamanla değişir'' kısımlarını $\checkmark$ ile işaretlemiş; $y(t) = t^2\cdot x(t-1)$ altı çizilmiş. Yeni el yazısı eklemeler:)}
\[
x(t-t_0) \longrightarrow y(t-t_0)
\]
\[
\begin{aligned}
x_2(t) &\longrightarrow y_2(t) = t^2\cdot x_2(t-1)\\
x_2(t) = x(t-t_0) &\qquad\ \ y_2(t) = t^2\cdot x(t-t_0-1)\\
&\qquad\ \ \underline{\underline{y(t-t_0)}} = (t-t_0)^2\,x(t-t_0-1)
\end{aligned}
\ \left.\vphantom{\begin{aligned}a\\a\end{aligned}}\right\}\neq
\]

\noindent\sayfa{41}
\sekilnotu[41]{Bu sayfa basılı materyal içeriyor (PDF görüntüleyicide açılmış çoktan seçmeli sınav sorusu, 4. soru, 10 puan: dürtü cevabı verilen ayrık zamanlı sistemin $u[n-3]$ girişi için çıkışı); dijitalleştirilmedi.}
El yazısı eklemeler:
\[
\kutula{x[n] * \delta[n-n_0] = x[n-n_0]}
\qquad,\qquad
x[n]\cdot\delta[n-n_0] = x[n_0]\,\delta[n-n_0]
\]
\begin{center}
\begin{tikzpicture}[>=Stealth]
  \node[draw,murekkep,thick,minimum width=1.4cm,minimum height=0.8cm] (h) at (0,0) {$h[n]$};
  \draw[->,murekkep,thick] (-2.2,0) node[above right] {$x[n]$} -- (h.west);
  \draw[->,murekkep,thick] (h.east) -- (1.6,0) node[right] {$y[n] = x[n] * h[n]$};
\end{tikzpicture}
\end{center}
\[
\begin{aligned}
y[n] &= \bigl(u[n-3]\bigr) * \bigl(-\delta[n] + \delta[n+2]\bigr)\\
&= -\underset{\uparrow}{\delta[n]} * \underset{\sim\sim\sim}{u[n-3]} + \delta[\underset{\sim\sim}{n+2}] * u[n-3]\\
&= -u[n-3] + u[n+2-3]\\
&= u[n-1] - u[n-3]
\end{aligned}
\]
\begin{center}
\begin{tikzpicture}[x=0.75cm,y=1cm]
  \draw[-{Stealth},murekkep] (-2.6,0) -- (6.2,0) node[right] {$n$};
  \fill[murekkep] (-2,0) circle (2pt) (-1,0) circle (2pt);
  \foreach \n in {1,2} {\draw[kirmizikalem,very thick] (\n,0) -- (\n,1); \fill[kirmizikalem] (\n,1) circle (2.5pt);
    \node[below,kirmizikalem] at (\n,0) {$\n$};}
  \foreach \n in {3,4,5} {\draw[murekkep,thick] (\n,0) -- (\n,1); \fill[murekkep] (\n,1) circle (2pt);}
  \node[below,murekkep] at (3,0) {$3$};
  \draw[murekkep,dashed] (5,1) -- (5.9,1);
  \node[murekkep,above right] at (5.2,1) {$u[n-1]$};
  \node[morkalem,right] at (6.4,0.35) {$u[n-3]$};
\end{tikzpicture}
\end{center}

\end{document}
